\documentclass[11pt]{article}

\usepackage[margin=1in]{geometry}
\usepackage{times}
\usepackage{microtype}
\usepackage{hyperref}
\usepackage{url}
\usepackage{graphicx}
\usepackage{booktabs}
\usepackage{array}
\usepackage{tabularx}
\usepackage{multirow}
\usepackage{xcolor}
\usepackage{listings}
\usepackage{caption}
\usepackage{subcaption}
\usepackage{amsmath}
\usepackage{parskip}
\usepackage{titlesec}
\usepackage{enumitem}
\usepackage{fancyhdr}
\usepackage{natbib}

\definecolor{codebg}{RGB}{248,248,248}
\definecolor{codeframe}{RGB}{200,200,200}
\definecolor{darkblue}{RGB}{26,35,126}
\definecolor{accentblue}{RGB}{41,121,255}

\lstset{
  basicstyle=\ttfamily\small,
  backgroundcolor=\color{codebg},
  frame=single,
  framesep=4pt,
  rulecolor=\color{codeframe},
  breaklines=true,
  breakatwhitespace=true,
  showstringspaces=false,
  captionpos=b,
  numbers=none,
}

\hypersetup{
  colorlinks=true,
  linkcolor=darkblue,
  citecolor=darkblue,
  urlcolor=accentblue,
}

\titleformat{\section}{\large\bfseries}{{\thesection.}}{0.5em}{}[\vspace{-0.3em}\rule{\linewidth}{0.4pt}]
\titleformat{\subsection}{\normalsize\bfseries}{{\thesubsection}}{0.5em}{}

\pagestyle{fancy}
\fancyhf{}
\rhead{\small MARS: Multi-Agent Research System}
\lhead{\small Shroff, 2026}
\rfoot{\thepage}

\setlength{\parskip}{0.6em}
\setlength{\parindent}{0pt}

% ──────────────────────────────────────────────────────────────────────────────
\begin{document}

\begin{center}
  {\LARGE\bfseries MARS: A Multi-Agent Research System with Cyclic\\[4pt]
  Quality Control and Chunk-Level Hallucination Grounding}

  \vspace{1.2em}
  {\large Harsh Shroff}\\[4pt]
  {\normalsize Independent Researcher\\
  \href{mailto:harshrofff@gmail.com}{\texttt{harshrofff@gmail.com}}}

  \vspace{0.6em}
  {\small March 2026}
\end{center}

\vspace{0.8em}
\hrule
\vspace{0.8em}

% ── ABSTRACT ─────────────────────────────────────────────────────────────────
\begin{abstract}
We present MARS (Multi-Agent Research System), an autonomous pipeline that
orchestrates 11 specialized LLM agents via a LangGraph state machine to produce
citation-grounded research reports on arbitrary topics. The system introduces
three principal contributions: (1)~a cyclic quality control loop in which a
dedicated Planner agent decomposes QC failure into structured search
directives---\texttt{search\_focus}, \texttt{suggested\_queries}, and
\texttt{tool\_priorities}---before re-routing to the Researcher, replacing the
generic retry behavior common in existing agentic pipelines; (2)~strict
chunk-level source grounding where every generated claim is bound to an indexed
\texttt{[SRC-X]} tag, enabling the Citation Verifier to detect and flag
referential errors by regex scan; and (3)~asynchronous parallel critic execution
via \texttt{asyncio.gather} across Gap Analyst, Bias Detector, and Fact Checker
agents, replacing sequential critique with true I/O concurrency. We evaluate
MARS across 15 runs spanning diverse topic domains and report a 0\%
\emph{citation tag error rate} (referential integrity), an 87\% QC first-pass
rate, and a mean cost of \$0.0028 per report at 130~seconds mean latency. We
explicitly distinguish citation tag integrity from semantic faithfulness, discuss
two QC failure modes arising in contested-evidence domains, quantify the
Human-in-the-Loop latency overhead (2.8$\times$), and characterize a Researcher
throughput ceiling as a bounded engineering limitation. Code is available at
\url{https://github.com/HarshShroff/multi-agent-researcher}.
\end{abstract}

\textbf{Keywords:} multi-agent systems, LangGraph, hallucination detection, quality
control, agentic AI, retrieval-augmented generation, human-in-the-loop

\hrule

% ── 1. INTRODUCTION ──────────────────────────────────────────────────────────
\section{Introduction}

The automation of literature synthesis and knowledge aggregation is a persistent
goal in information retrieval. Recent advances in large language models (LLMs)
have produced a class of agentic pipelines in which multiple specialized LLM
calls are composed to decompose and solve complex tasks. Yet the dominant failure
mode of such systems is hallucination: the generation of plausible but
unsupported claims \citep{ji2023hallucination}. In retrieval-augmented
generation (RAG) settings, this takes the specific form of claims that cite real
sources but misrepresent their content---a failure that structural citation
enforcement alone cannot catch.

Existing multi-agent research tools---including AutoGPT \citep{toran2023autogpt},
BabyAGI \citep{nakajima2023babyagi}, and LangChain's plan-and-execute
patterns---address task decomposition but lack principled mechanisms for QC
failure analysis and targeted remediation. When quality control is present, it
typically takes the form of a single-pass LLM self-evaluation, which provides no
directional signal on failure.

This paper describes MARS (Multi-Agent Research System), designed to address
these gaps through three mechanisms. First, we replace generic retry logic with a
Planner-mediated cyclic loop: on QC failure, a dedicated Planner agent analyzes
the failure and emits a typed JSON strategy object that directs the next
Researcher iteration toward the specific coverage gap. Second, we enforce
chunk-level source grounding---every claim the Writer produces must embed an
\texttt{[SRC-X]} tag referencing a specific indexed chunk, and the Citation
Verifier flags any tag pointing to a non-existent chunk. Third, we implement a
toggleable Human-in-the-Loop (HITL) pause via LangGraph's
\texttt{interrupt\_before} mechanism, with explicit measurement of its latency
and token cost.

A point of precision: the 0\% citation tag error rate we report measures
\emph{referential integrity}---whether \texttt{[SRC-X]} tags point to real
indexed chunks. It does not measure \emph{semantic faithfulness}---whether the
claim adjacent to a tag is actually supported by that chunk's content. Full
faithfulness evaluation via RAGAS \citep{es2023ragas} or LLM-as-judge is
deferred to future work. We state this distinction explicitly because conflating
the two has been a documented source of misleading RAG evaluation claims.

We evaluate MARS across 15 controlled runs spanning eight topic domains and four
report depth settings, logging QC scores, citation tag errors, retry counts,
per-agent token costs, and source yield per run.

\vspace{0.5em}
The paper is organized as follows: Section~2 reviews related work. Section~3
describes the system architecture. Section~4 details agent design. Section~5
presents evaluation results. Section~6 discusses findings and limitations.
Section~7 concludes.

% ── 2. RELATED WORK ──────────────────────────────────────────────────────────
\section{Related Work}

\subsection{Multi-Agent LLM Systems}

The multi-agent paradigm was introduced to LLM pipelines by AutoGPT
\citep{toran2023autogpt} and BabyAGI \citep{nakajima2023babyagi}, which showed
that iterative self-prompting loops could tackle open-ended tasks. LangChain
\citep{chase2022langchain}, CrewAI \citep{moura2024crewai}, and Microsoft's
AutoGen \citep{wu2023autogen} provided infrastructure for role-specialized agents
and inter-agent communication. LangGraph \citep{langgraph2024} extended this by
modeling agent pipelines as directed graphs with state persistence, conditional
edges, and interrupt hooks---the substrate on which MARS is built.

The key architectural distinction from prior work: MARS does not decompose a task
into subtasks and aggregate their outputs in a single pass. It implements a
stateful retry loop with typed failure analysis, treating QC failure as a signal
carrying information about coverage gaps rather than as a trigger for generic
re-execution.

\subsection{Hallucination in RAG Pipelines}

Retrieval-augmented generation \citep{lewis2020rag} reduces hallucination by
grounding generation in retrieved documents, but does not eliminate it.
\citet{shi2023irrelevant} showed that LLMs can fail to use retrieved context even
when it is present, generating from parametric memory instead. RAGAS
\citep{es2023ragas} introduced reference-free faithfulness evaluation by checking
whether generated answers are entailed by retrieved chunks at the sentence level.

Our chunk-level grounding approach differs in mechanism from RAGAS: rather than
post-hoc faithfulness scoring, we enforce grounding structurally by requiring
every claim to embed a verifiable chunk tag, detecting referential violations at
generation time. This eliminates phantom citation failures---tags pointing to
non-existent sources---but does not address semantic faithfulness, which requires
content-level comparison between claim and chunk.

\subsection{Quality Control in Agentic Systems}

Self-refinement \citep{madaan2023selfrefine} showed that LLMs can improve their
outputs by critiquing and revising their own generations, but this conflates
generation and evaluation in a single model, limiting critique independence.
Reflexion \citep{shinn2023reflexion} introduced verbal reinforcement via episodic
memory.

MARS differs from both: QC is performed by a separate agent with no generative
stake in the research, and on failure, a third agent (the Planner) independently
generates a search remediation strategy. This maintains separation of concerns
across generation, evaluation, and planning---and avoids the sycophancy risk that
arises when a model evaluates its own outputs \citep{sycophancy2023}.

\subsection{Human-in-the-Loop Oversight}

\citet{seshia2018formal} and \citet{anthropic2023modelspec} highlight the
importance of human oversight in autonomous AI systems. LangGraph's
\texttt{interrupt\_before} primitive provides a natural HITL integration point.
Prior systems such as HuggingGPT \citep{shen2023hugginggpt} include human
verification steps, but few have explicitly measured the latency and token cost
of human intervention. We provide this measurement in Section~5.4.

% ── 3. SYSTEM ARCHITECTURE ───────────────────────────────────────────────────
\section{System Architecture}

MARS is implemented as a directed cyclic graph using LangGraph's
\texttt{StateGraph} API. The state schema (\texttt{ResearchState}) is a
\texttt{TypedDict} carrying all inter-agent data: the research topic,
configuration, retrieved chunks, analysis outputs, critic outputs, QC score and
feedback, planner strategy, and an append-only log list.

\subsection{Pipeline Overview}

\begin{figure}[h]
\centering
\begin{lstlisting}[basicstyle=\ttfamily\footnotesize, backgroundcolor=\color{codebg}]
User Input & PDFs
      |
      v
[Researcher: ReAct + MCP Tools]  <-- planner_strategy on retry
      |  indexed chunks [SRC-0..N]
      v
[Analyst: Pydantic structured output]
      |
      v
Parallel Critics (asyncio.gather):
  [Gap Analyst] | [Bias Detector] | [Fact Checker]
      |
      v
[Synthesizer]
      |
      v
[QC Agent]  -- score 0-100, threshold 85
      |
  FAIL |--------------------------------> [Planner]
  (qc_passed=False)                            |
                                  {search_focus, suggested_queries,
                                   tool_priorities}  --> Researcher
  PASS |
      v
[HITL Gate]  (toggleable via interrupt_before)
      |  approve / inject feedback
      v
[Formatter] --> [Writer: streaming] --> [Citation Verifier]
                                              |
                                              v
                              Final Report + Visualizations
                              PDF | Markdown | JSON | TXT
\end{lstlisting}
\caption{MARS agent pipeline. The cyclic path through the Planner is the core architectural contribution.}
\label{fig:pipeline}
\end{figure}

\subsection{State Machine and Graph Edges}

The graph is compiled with a \texttt{MemorySaver} checkpointer, enabling full
state persistence per thread and supporting \texttt{interrupt\_before} for HITL.
A conditional edge after the QC node evaluates \texttt{qc\_passed} from the
state and routes to the Planner (failure) or the HITL gate (success). A
secondary conditional edge after the HITL gate routes to the Formatter on user
approval, or back to the Researcher on injected feedback.

The cyclic Researcher $\to \cdots \to$ QC $\to$ Planner $\to$ Researcher
structure is the central architectural novelty. Unlike systems that terminate on
first QC pass or retry generically, MARS threads the Planner's typed strategy
through the state dict to guide each subsequent Researcher iteration, preventing
degenerate re-execution of the same queries.

\subsection{MCP Tool Decoupling}

Search tools (ArXiv, DuckDuckGo, Wikipedia) are exposed via a standalone MCP
server (\texttt{mcp\_server.py}) over stdio. The Researcher communicates with
this server via the Model Context Protocol \citep{mcp2024} rather than invoking
tool functions directly. This decoupling makes the data-gathering pipeline
platform-agnostic: the same MCP server can serve any MCP-compliant client
without modification to the tool implementations.

% ── 4. AGENT DESIGN ──────────────────────────────────────────────────────────
\section{Agent Design}

All agents inherit from a shared \texttt{Agent} base class wrapping Gemini~2.5
Flash Preview (\texttt{gemini-2.5-flash-preview}) via the
\texttt{google-generativeai} SDK. The base class provides \texttt{call\_llm()}
and \texttt{call\_llm\_async()} methods, each updating a thread-local token
counter for per-agent cost attribution.

\subsection{Researcher Agent (ReAct + MCP)}

The Researcher executes a ReAct loop \citep{yao2022react}, interleaving reasoning
steps with tool calls to the MCP server. On each iteration, it receives the
current \texttt{planner\_strategy} (if any) and uses \texttt{suggested\_queries}
and \texttt{tool\_priorities} to construct targeted queries rather than generic
topic search. Source deduplication is performed by normalizing each URL
(stripping query parameters and fragments) before indexing, preventing redundant
chunks from repeated tool calls. Each retrieved document is parsed into a chunk
with fields \texttt{content}, \texttt{source\_url}, \texttt{source\_name}, and
\texttt{source\_type} (ArXiv / Web / Wikipedia / Uploaded PDF), indexed by a
sequential \texttt{[SRC-X]} key in the state's \texttt{research\_chunks} dict.

\subsection{Quality Control Agent}

The QC agent receives the topic, a \texttt{research\_data} summary, analysis
output, and \texttt{qc\_iterations} count. It scores research on a 0--100 scale
via a structured prompt penalizing missing sub-topics, thin source coverage, and
internal contradictions. Scores below 85 trigger \texttt{qc\_passed: false} with
\texttt{qc\_feedback} listing identified deficiencies.

An important methodological note: the QC agent uses the same foundation model
(Gemini Flash) that generated the research. This introduces a self-evaluation
bias documented in the LLM sycophancy literature \citep{sycophancy2023}. The QC
scores we report reflect model self-assessment, not ground-truth quality
measurement. Using an independent, stronger judge model is listed as future work.

The 85-point threshold was chosen empirically to achieve a first-pass rate above
80\% across our evaluation set while still triggering retries on demonstrably
thin research. At this setting the system achieves an 87\% first-pass rate.

\subsection{Planner Agent}

The Planner is activated exclusively on QC failure. It receives
\texttt{qc\_feedback}, the list of previously retrieved sources, and the topic,
and outputs a JSON object conforming to a Pydantic schema:

\begin{lstlisting}[language=Python]
class PlannerStrategy(BaseModel):
    search_focus: str           # coverage gap, one sentence
    suggested_queries: list[str]   # 3-5 targeted queries
    tool_priorities: list[str]     # ['arxiv', 'web', 'wiki'] ordered
\end{lstlisting}

Explicit structuring of the remediation---rather than passing QC feedback
verbatim to the Researcher---forces the Planner to synthesize an actionable
directive. This is analogous to the distinction between random restart and
gradient descent in optimization: the Planner provides a directional signal
rather than returning the Researcher to its original search state.

A limitation of this mechanism: Planner strategy quality depends on QC feedback
quality, which in turn depends on the QC agent's ability to identify gaps from a
summary view. For epistemological gaps (contested evidence, rapidly evolving
benchmarks) rather than topical gaps (missing subtopics), the Planner may
generate topically valid but substantively insufficient queries, as observed in
Runs~2 and~10 (Section~5.3).

\subsection{Citation Verifier}

After the Writer produces the draft, the Citation Verifier performs two
operations. First, it extracts all \texttt{[SRC-X]} tags via regex and
cross-references each index against \texttt{research\_chunks}. Tags referencing
non-existent indices are flagged as referential errors and counted. The citation
tag error rate is computed as:

\[
\text{tag error rate} = \frac{\text{invalid tags}}{\text{valid tags} + \text{invalid tags}}
\]

Second, it replaces valid tags with formatted inline citations (IEEE, APA~7th,
MLA~9th, Harvard, or Chicago) and appends a formatted references section.

We reiterate the scope of this metric: it measures whether the LLM generated
syntactically valid citation tags pointing to real chunks. It does not measure
whether the claims adjacent to those tags are actually supported by the chunk
content. A claim that misrepresents a real source will not be caught by this
verifier. Semantic faithfulness evaluation via RAGAS or an LLM-as-judge is
deferred to future work.

\subsection{Async Parallel Critics}

The Gap Analyst, Bias Detector, and Fact Checker are dispatched concurrently via
\texttt{asyncio.gather} inside a LangGraph async node. Each receives a snapshot
of \texttt{research\_data} and emits a structured critique string. The
Synthesizer merges these three strings into the context window passed to QC.
Concurrency is achieved via \texttt{call\_llm\_async()}, which releases the
event loop during network I/O, reducing wall time relative to sequential critic
execution.

% ── 5. EVALUATION ────────────────────────────────────────────────────────────
\section{Evaluation}

\subsection{Experimental Protocol}

We conducted 15 runs across eight topic domains (mechanistic interpretability,
materials science, computer vision, neuromorphic computing, biomedical AI,
robotics, sustainability, AI policy), four report depth settings (Concise,
Standard, Detailed, Exhaustive), and two HITL configurations. All runs used
\texttt{gemini-2.5-flash-preview} as the backbone model. A custom
\texttt{experiment\_logger} module appended one JSON record per run to a JSONL
log, capturing: \texttt{run\_id}, timestamp, topic, depth, HITL flag, QC events
(score and pass/fail per iteration), \texttt{total\_retries}, \texttt{wall\_time\_s},
per-agent token breakdown, grounding statistics, and per-source-type chunk
counts. All 15 runs were executed on March~9, 2026. Runs~13 and~14 are matched
pairs on the identical topic (``Agentic AI: autonomy, safety, and deployment'')
with HITL enabled and disabled respectively to enable direct comparison.

\subsection{Main Results}

\begin{table}[h]
\centering
\caption{Aggregate metrics across 15 runs.}
\label{tab:aggregate}
\begin{tabular}{lr}
\toprule
\textbf{Metric} & \textbf{Value} \\
\midrule
Total runs & 15 \\
QC first-pass rate & 87\% (13/15) \\
Mean QC score (passing runs) & 91.5 / 100 \\
QC failures (unresolved in 2 retries) & 2 (Runs 2, 10) \\
Mean retries per run & 0.27 \\
Citation tag error rate & 0.00\% \\
Mean wall time & 130.2 s \\
Mean cost per report & \$0.0028 \\
Mean total tokens per run & 30,874 \\
Mean unique sources cited & 10.5 \\
Mean total chunks indexed & 10.7 \\
\bottomrule
\end{tabular}
\end{table}

\begin{table}[h]
\centering
\caption{Per-run results. Asterisk (*) denotes final score after Planner-mediated retry; threshold not reached within 2 iterations.}
\label{tab:runs}
\small
\begin{tabular}{clllrrrr}
\toprule
\textbf{Run} & \textbf{Topic (abbreviated)} & \textbf{Depth} & \textbf{QC} & \textbf{Pass} & \textbf{Retries} & \textbf{Time (s)} & \textbf{Cost (\$)} \\
\midrule
0  & Mechanistic interpretability     & Detailed   & 92    & \checkmark & 0 & 237 & 0.0029 \\
1  & Solid-state batteries            & Standard   & 95    & \checkmark & 0 &  63 & 0.0020 \\
2  & Diffusion models vs.\ GANs       & Detailed   & 60*   & $\times$   & 2 & 243 & 0.0030 \\
3  & Neuromorphic computing           & Standard   & 89    & \checkmark & 0 & 102 & 0.0029 \\
4  & CRISPR base editing              & Detailed   & 97    & \checkmark & 0 & 115 & 0.0030 \\
5  & Humanoid robotics market         & Standard   & 94    & \checkmark & 0 & 125 & 0.0031 \\
6  & OSS vs.\ proprietary LLMs        & Standard   & 91    & \checkmark & 0 & 119 & 0.0030 \\
7  & Carbon capture economics         & Concise    & 88    & \checkmark & 0 & 113 & 0.0029 \\
8  & Consciousness in AI systems      & Detailed   & 94    & \checkmark & 0 & 120 & 0.0032 \\
9  & AGI timelines                    & Standard   & 88    & \checkmark & 0 & 123 & 0.0033 \\
10 & Psychedelic-assisted therapy     & Standard   & 65*   & $\times$   & 2 & 164 & 0.0035 \\
11 & Transformer architecture         & Exhaustive & 91    & \checkmark & 0 & 109 & 0.0030 \\
12 & Multi-agent AI systems           & Exhaustive & 91    & \checkmark & 0 & 131 & 0.0031 \\
13 & Agentic AI (HITL ON)             & Standard   & 96    & \checkmark & 0 & 153 & 0.0027 \\
14 & Agentic AI (HITL OFF)            & Standard   & 94    & \checkmark & 0 &  54 & 0.0014 \\
\bottomrule
\end{tabular}
\end{table}

\subsection{QC Failure Mode Analysis}

Both QC failures (Runs~2 and~10) involve topics with contested empirical claims
and heterogeneous source quality. Run~2 (Diffusion models vs.\ GANs) involves a
rapidly evolving benchmark landscape where ArXiv papers disagree on key metrics
depending on publication date. Run~10 (Psychedelic-assisted therapy) involves a
domain where regulatory status, clinical trial heterogeneity, and secondary
coverage introduce substantial noise relative to primary RCT evidence.

In both cases, QC feedback cited ``insufficient coverage of counterarguments''
and ``over-reliance on secondary sources.'' The Planner responded with targeted
queries directing the Researcher toward primary literature. Despite two retry
iterations, neither run reached the 85-point threshold.

We interpret this as graceful degradation rather than system failure. The QC
mechanism correctly identified that the available evidence base did not support a
high-confidence synthesis, and surfaced this explicitly rather than producing an
overconfident report. This behavior is preferable to a system that passes QC on
contested topics by generating plausible-sounding but unsupported claims.

\subsection{HITL vs.\ Automated Latency}

Runs~13 and~14 are matched pairs on identical topics. HITL-enabled wall time was
153~s vs.\ 54~s for the automated run (2.8$\times$ overhead). Cost was \$0.0027
vs.\ \$0.0014 (1.9$\times$ overhead).

The wall-time overhead includes human review time, making 2.8$\times$ an upper
bound on HITL protocol cost. The token cost overhead (1.9$\times$) reflects a
second Writer invocation: in HITL mode, the Writer generates a pre-approval
draft for review before producing the final version, approximately doubling
Writer token expenditure. Separating computational overhead from human review
time requires logging \texttt{interrupt\_before} entry and exit timestamps
separately, which we defer as a future instrumentation change.

\subsection{Source Yield and Depth Analysis}

Chunk counts plateau at 10--12 across all 15 runs regardless of depth setting,
with a breakdown of approximately 5 ArXiv + 5 web + 0--2 Wikipedia chunks per
run. This uniformity reflects the Researcher's per-tool cap of
\texttt{max\_results=5}. The depth setting controls only the Writer's context
window allocation (30~/ 60~/ 120~/ 200 chunks for Concise through Exhaustive),
not retrieval volume.

The consequence: Exhaustive mode directs the Writer to produce a longer report
from the same data, which risks verbose elaboration and potential semantic drift
at high word counts. This interaction is not tracked in the current evaluation
and represents the highest-risk configuration for semantic faithfulness failures.

% ── 6. DISCUSSION ────────────────────────────────────────────────────────────
\section{Discussion}

\subsection{Structured Planner vs.\ Generic Retry}

The Planner-mediated retry is the core architectural contribution of MARS. In
systems routing QC failures directly to the Researcher, the Researcher has no
additional information about why prior research was insufficient; it may retrieve
the same sources and revisit the same queries. The Planner breaks this degeneracy
by synthesizing QC feedback into a typed JSON directive. The analogy to gradient
descent vs.\ random restart is apt: the Planner provides directional information
that guides the Researcher toward unexplored coverage regions.

We note a limitation of this mechanism not yet addressed experimentally: we have
not run a controlled ablation comparing Planner-directed retry against
raw-feedback retry. Until we do, the claim that structured strategy improves over
direct QC feedback pass-through is an architectural hypothesis rather than an
empirical result. This ablation is listed as a priority for future work.

\subsection{Limitations}

\begin{enumerate}[leftmargin=*, label=\textbf{L\arabic*.}]

\item \textbf{Semantic faithfulness not measured.} The 0\% citation tag error
rate measures referential integrity only. Semantic faithfulness evaluation via
RAGAS \citep{es2023ragas} or LLM-as-judge \citep{zheng2023judging} is required
to support stronger claims about hallucination.

\item \textbf{Self-evaluation bias.} The QC agent uses the same model that
generated the research. Mean QC scores of 91.5 reflect self-assessment and
should not be interpreted as ground-truth quality measurements without
validation against a stronger independent judge.

\item \textbf{Researcher throughput ceiling.} The per-tool cap of
\texttt{max\_results=5} limits chunk yield to approximately 10--15 per run
regardless of depth. Depth-responsive retrieval scaling is a straightforward
engineering change deferred as future work.

\item \textbf{HITL measurement artifact.} Reported wall-time overhead includes
human review time. The pure computational overhead is approximately 1.9$\times$
in tokens, reflecting the second Writer invocation.

\item \textbf{Evaluation breadth.} N=15 across 8 topic domains does not support
statistical inference. Results are directional. A larger evaluation set with
controlled topic sampling is needed for benchmark-quality claims.

\item \textbf{Planner ablation missing.} The claim that structured Planner
strategy outperforms raw QC feedback retry is not yet empirically validated.

\end{enumerate}

\subsection{Broader Implications}

MARS shows that chunk-level citation grounding and structured QC remediation are
achievable at sub-cent cost using current frontier model APIs. The MCP-based
tool decoupling suggests a path toward interoperable agent pipelines composable
with other MCP-compliant tools---code executors, database connectors, enterprise
APIs---without architectural modification. The patterns introduced here (typed
Planner remediation, chunk-indexed grounding, async parallel critics) are
composable primitives applicable beyond text synthesis to other agentic research
tasks.

% ── 7. CONCLUSION ────────────────────────────────────────────────────────────
\section{Conclusion}

We presented MARS, a multi-agent research system built on LangGraph that
introduces Planner-mediated cyclic QC with structured search remediation,
chunk-level citation tag grounding with runtime referential error detection, and
asynchronous parallel critic execution. Across 15 runs, MARS achieves a 0\%
citation tag error rate, an 87\% QC first-pass rate, and a mean cost of \$0.0028
per report.

We have been explicit throughout about what these metrics do and do not measure.
The citation tag error rate catches phantom citations but not semantic
misrepresentation. The QC scores reflect self-assessment under known sycophancy
conditions. The N=15 evaluation supports directional conclusions, not benchmarks.
These limitations are engineering problems with clear paths to resolution, not
structural flaws in the approach.

Future work includes depth-responsive retrieval, a Planner ablation study,
semantic faithfulness evaluation via RAGAS, and an independent judge model for
QC validation. The system is open source at
\url{https://github.com/HarshShroff/multi-agent-researcher}.

% ── REFERENCES ───────────────────────────────────────────────────────────────
\bibliographystyle{plainnat}
\bibliography{references}

\end{document}
